% Generated draft: requires mathematical and source review.
\documentclass[11pt]{article}
\usepackage[margin=25mm]{geometry}
\usepackage{fontspec}
\setmainfont{Latin Modern Roman}
\usepackage{amsmath,amssymb,mathtools,unicode-math}
\setmathfont{Latin Modern Math}
\usepackage{enumitem,needspace,xurl,hyperref}
\setlength{\parindent}{0pt}
\setlength{\parskip}{5pt}
\setlength{\emergencystretch}{3em}
\newcommand{\N}{\mathbb{N}}
\newcommand{\Z}{\mathbb{Z}}
\newcommand{\Q}{\mathbb{Q}}
\newcommand{\R}{\mathbb{R}}
\newcommand{\C}{\mathbb{C}}
\newcommand{\F}{\mathbb{F}}
\newcommand{\A}{\mathbb{A}}
\newcommand{\PP}{\mathbb{P}}
\newcommand{\E}{\mathbb{E}}
\newcommand{\eps}{\varepsilon}
\newcommand{\dd}{\,\mathrm{d}}
\begin{document}
\section*{UnsolvedMath ID: 700761}
\textbf{Source problem code:} catalogue entry
\subsection{Definitions and mathematical statement}
Let \(\Omega\) be an infinite set and let \(G\leq \operatorname{Sym}(\Omega)\) be a permutation group. Define
\[
\Omega^{1,\{2\}}=\{(\alpha,\{\beta,\gamma\}):\alpha,\beta,\gamma\in\Omega\text{ are pairwise distinct}\}.
\]
The induced action is
\[
g\cdot(\alpha,\{\beta,\gamma\})=(g\alpha,\{g\beta,g\gamma\}).
\]
Call \(G\) \(1\text{-}\{2\}\)-transitive if this action is transitive. For \(\alpha\in\Omega\), write \(G_\alpha=\{g\in G:g\alpha=\alpha\}\). A group acting on a set \(X\) is \(2\)-homogeneous if it is transitive on the unordered two-element subsets of \(X\). Thus \(G\) is \(1\text{-}\{2\}\)-transitive if and only if it is transitive on \(\Omega\) and \(G_\alpha\) is \(2\)-homogeneous on \(\Omega\setminus\{\alpha\}\) for every \(\alpha\in\Omega\).

Write
\[
\Omega^{(3)}=\{(\alpha,\beta,\gamma)\in\Omega^3:\alpha,\beta,\gamma\text{ are pairwise distinct}\}.
\]
The group is \(3\)-transitive if its coordinatewise action on \(\Omega^{(3)}\) is transitive.

\textbf{Classification problem (P. M. Neumann, 1990).} Classify all infinite permutation groups \(G\leq\operatorname{Sym}(\Omega)\) that are \(1\text{-}\{2\}\)-transitive but are not \(3\)-transitive.

The source places ``infinite'' in parentheses; the formulation above states the infinite case. Within this transitive permutation-group setting, infinitude of \(G\) and infinitude of \(\Omega\) are equivalent. Whether the parenthetical wording also requests a finite-case classification is not explicit. The source also does not specify the equivalence relation under which the classification is sought.
\subsection{Short English statement}
Classify the infinite permutation groups that can carry any distinguished point together with an unordered pair of other points to any other such configuration, but cannot carry every ordered triple of distinct points to every other ordered triple. The scope of the parenthetical word ``infinite'' and the intended equivalence for classification require clarification.
\subsection{Sources}
\begin{itemize}
\item Local UnsolvedMath registry, supplied complete record with permanent ID \(700761\), exact source code \texttt{catalogue entry}, title \emph{Kourovka Notebook Problem 11.69}; upstream revision \texttt{d371ee917ac9647fbb06cd6fc529deec6f3dda82}.
\item \emph{The Kourovka Notebook (Unsolved Problems in Group Theory)}, Problem 11.69, 11th issue (1990), proposed by P. M. Neumann. The supplied background also identifies the 21st issue edition (2026), arXiv:1401.0300v48. The supplied collection description names E. I. Khukhro and V. D. Mazurov as editors.
\item The supplied registry labels this problem \texttt{open} and records an update date of 2026-10-16. This is supplied registry metadata, not an independently verified assessment of the literature or current mathematical status.
\end{itemize}
\end{document}
